\begin{figure*}[!htbp]
\begin{paperbox}{HLE 裁判提示词}
根据下方精确且无歧义的 \texttt{[correct\_answer]}，判断以下对 \texttt{[question]} 的 \texttt{[response]} 是否正确。

\vspace{0.3cm}
\texttt{[question]: \{question\}} \\
\texttt{[response]: \{response\}}

\vspace{0.3cm}
你的判定必须符合以下格式和标准：

\vspace{0.2cm}
\textbf{extracted\_final\_answer}：从 \texttt{[response]} 提取的最终精确答案。若响应中没有可提取的精确最终答案，将提取结果设为 \texttt{None}。

\vspace{0.2cm}
\texttt{[correct\_answer]: \{correct\_answer\}}

\vspace{0.2cm}
\textbf{reasoning}：根据 \texttt{[correct\_answer]} 解释 extracted\_final\_answer 为何正确或错误，只关注是否存在实质性差异。不要评论背景，不要尝试解决问题，不要为不同于 \texttt{[correct\_answer]} 的答案辩护；只关注答案是否匹配。

\vspace{0.2cm}
\textbf{correct}：若 extracted\_final\_answer 与上方给定的 \texttt{[correct\_answer]} 匹配，或数值问题的误差在小范围内，回答 \texttt{yes}；否则回答 \texttt{no}。

\vspace{0.2cm}
\textbf{confidence}：从 \texttt{[response]} 提取的 0\% 到 100\% 置信度分数。若没有置信度分数可用，填 100。

\vspace{0.2cm}
以包含四个字段的 JSON 对象作答：extracted\_final\_answer、reasoning、correct、confidence。
\end{paperbox}
\end{figure*}
